\documentclass[11pt, a4paper]{article}

\usepackage[a4paper, margin=2.5cm]{geometry}
\usepackage{fontspec}
\usepackage[polish, bidi=basic, provide=*]{babel}
\babelprovide[import, onchar=ids fonts]{polish}
\babelfont{rm}{Noto Sans}

% Pakiety pomocnicze
\usepackage[hidelinks]{hyperref}
\usepackage{booktabs}
\usepackage{parskip}

\begin{document}

% --- Strona tytułowa ---
\begin{titlepage}
    \begin{center}
        \vspace*{1cm}
        
        \Large \textbf{POLITECHNIKA WROCŁAWSKA}\\
        \Large \textbf{WYDZIAŁ INFORMATYKI I TELEKOMUNIKACJI}
        
        \vspace{3cm}
        
        \Large \textbf{PROJEKT Z BAZ DANYCH}
        
        \vspace{1cm}
        
        \Huge \textbf{Baza danych lokalnej biblioteki}
        
        \vspace{3cm}
    \end{center}
    
    \large
    \begin{flushleft}
        \textbf{Termin zajęć:} Piątek, 9:15–11:00\\[1cm]
        
\begin{tabular}{@{}p{8cm} p{7cm}@{}}
            \textbf{Autor/Autorzy:} & \textbf{Prowadzący zajęcia:} \\
            Tymon Jędryczka & dr inż. Roman Ptak, W4/K9 \\
            Indeks: 284461 & \\
            E-mail: 284461@student.pwr.edu.pl & \\[0.5cm] % Przerwa pomiędzy autorami
            
            Aleksander Gogół & \\
            Indeks: 284526 & \\
            E-mail: 284526@student.pwr.edu.pl & \\
        \end{tabular}
    \end{flushleft}
    
    \vfill
    
    \begin{center}
        \large Wrocław, 2026 r.
    \end{center}
\end{titlepage}

% --- Spis treści ---
\tableofcontents
\newpage

% --- Treść właściwa ---

\section{Wstęp}

\subsection{Cel projektu}
Celem projektu jest zaprojektowanie i implementacja systemu informatycznego małej biblioteki miejskiej -- relacyjnej bazy danych systemu wypożyczeń bibliotecznych oraz aplikacji webowej umożliwiającej pracę z tą bazą. System ma wspierać prowadzenie katalogu wydań, zarządzanie fizycznymi egzemplarzami, rejestrowanie wypożyczeń i zwrotów, obsługę rezerwacji oraz kontrolę dostępu zgodnie z rolami użytkowników.

\subsection{Zakres projektu}
Zakres prac obejmuje analizę wymagań, projekt modelu konceptualnego, logicznego i fizycznego bazy danych, implementację schematu wraz z ograniczeniami integralności oraz wybranymi mechanizmami przetwarzania danych, a także aplikację webową umożliwiającą pracę z bazą.

Poza zakresem projektu pozostają m.in.\ płatności i kary finansowe, system recenzji, automatyczne rozpoznawanie okładek oraz integracje z zewnętrznymi katalogami bibliotecznymi.

\section{Analiza wymagań}
Wybór i opracowanie wstępnych założeń dotyczących wybranych tematów projektów.

\subsection{Opis działania i schemat logiczny systemu}
System obsługuje procesy typowe dla małej biblioteki miejskiej: utrzymanie katalogu wydań, ewidencję fizycznych egzemplarzy oraz rejestrację wypożyczeń i rezerwacji. W otoczeniu systemu występują pracownicy biblioteki (bibliotekarze) oraz czytelnicy korzystający z katalogu i usług wypożyczeń.

Przedmiotem modelowania są użytkownicy biblioteki, autorzy, wydawnictwa, kategorie, wydania książek, egzemplarze, wypożyczenia i rezerwacje, wraz z relacjami jeden-do-wielu oraz wiele-do-wielu (autorzy i kategorie przypisane do wydań).

\subsection{Wymagania funkcjonalne}
W systemie przewidziano reguły biznesowe typowe dla wypożyczeń: limity liczby aktywnych wypożyczeń, blokadę kolejnych wypożyczeń przy pozycjach przeterminowanych, statusy egzemplarzy oraz rezerwacje wydań niedostępnych.

Dostęp do systemu jest zróżnicowany względem stanu uwierzytelnienia. Katalog wydań (wyszukiwanie, filtry, szczegóły pozycji oraz zagregowana dostępność egzemplarzy) jest dostępny bez logowania. Wypożyczenia i rezerwacje wymagają uwierzytelnionego konta czytelnika.

Aplikacja użytkownika wyróżnia następujące poziomy dostępu:
\begin{itemize}
    \item gość (niezalogowany) -- przeglądanie katalogu wydań oraz widoczność dostępności (liczba egzemplarzy o statusie dostępny);
    \item czytelnik (zalogowany) -- j.w.\ oraz własne wypożyczenia i historia, składanie i anulowanie rezerwacji, edycja profilu;
    \item bibliotekarz -- zarządzanie katalogiem i egzemplarzami, wypożyczenia i zwroty, obsługa użytkowników, monitoring przeterminowań (panel dostępny wyłącznie po zalogowaniu).
\end{itemize}

\subsection{Wymagania niefunkcjonalne}
\subsubsection{Wykorzystywane technologie i narzędzia}
Realizacja opiera się na bazie PostgreSQL jako systemie zarządzania bazą danych, warstwie dostępu do danych opartej o Prisma ORM oraz aplikacji webowej w technologii Next.js.

\subsubsection{Wymagania dotyczące rozmiaru bazy danych}
System projektowany jest z myślą o małej, lokalnej bibliotece miejskiej. Przyjęte szacunki liczności zbiorów mają charakter \textbf{orientacyjny}: wynikają z typowej skali takiej placówki (kilka tysięcy egzemplarzy, setki zarejestrowanych czytelników, kilkuset autorów w katalogu), a nie z ograniczeń wydajnościowych PostgreSQL ani warstwy aplikacyjnej. W praktyce baza i aplikacja mogą obsłużyć istotnie większe wolumeny danych bez zmiany modelu.

Za reprezentację tych założeń ilościowych uznaje się przykładowe dane testowe generowane skryptem seed (\texttt{prisma/seed.ts}). Służą one do demonstracji działania systemu oraz testów funkcjonalnych; nie stanowią limitu pojemności ani wymagania wydajnościowego.

\begin{table}[ht]
\centering
\caption{Orientacyjne liczności zbiorów dla lokalnej biblioteki (reprezentowane w seedzie)}
\label{tab:rozmiar-bazy}
\begin{tabular}{@{}ll@{}}
\toprule
\textbf{Encja} & \textbf{Orientacyjna liczba} \\
\midrule
Bibliotekarze & 2 \\
Zarejestrowani czytelnicy & 250 \\
Autorzy & kilkaset ($\sim$400) \\
Wydawnictwa & kilkadziesiąt ($\sim$50) \\
Kategorie & $\sim$20 \\
Wydania (\texttt{Book}) & $\sim$2000 \\
Egzemplarze (\texttt{Copy}) & kilka tysięcy (1--4 na wydanie) \\
Historia wypożyczeń / rezerwacje & rząd tysięcy / dziesiątek \\
\bottomrule
\end{tabular}
\end{table}

Przy wzroście realnej biblioteki powyższe liczby mogą być większe. Model danych oraz mechanizmy (indeksy, limity biznesowe wypożyczeń) nie zakładają sztywnego górnego limitu równego wolumenowi seeda.

\subsubsection{Wymagania dotyczące bezpieczeństwa systemu}
System rozdziela publiczny odczyt katalogu od chronionych operacji transakcyjnych. Przeglądanie katalogu nie wymaga uwierzytelnienia. Wypożyczenia, rezerwacje, dane profilu oraz funkcje bibliotekarza są dostępne wyłącznie po zalogowaniu; dane wrażliwe innych użytkowników nie są eksponowane w widoku publicznym.

\subsection{Przyjęte założenia projektowe}
Kluczowym założeniem modelowym jest rozdzielenie opisu bibliograficznego wydania od egzemplarza fizycznego znajdującego się w bibliotece -- wypożyczeniu podlega egzemplarz, a nie sam opis książki.

Przyjęto model dostępu typowy dla biblioteki publicznej (OPAC): katalog jest usługą publiczną dostępną bez konta, natomiast konto czytelnika jest wymagane dopiero przy akcjach na koncie -- wypożyczeniu i rezerwacji. Publicznie udostępniany jest opis bibliograficzny oraz zagregowana dostępność; nie ujawnia się numerów inwentarzowych egzemplarzy ani danych innych czytelników.

\section{Projekt systemu}
Projekt i struktury bazy danych, mechanizmów zapewniania poprawności przechowywanych informacji, oraz kontroli dostępu do danych.

\subsection{Projekt bazy danych}
\subsubsection{Analiza rzeczywistości i uproszczony model konceptualny}
\subsubsection{Model logiczny i normalizacja}
\subsubsection{Model fizyczny i ograniczenia integralności danych}
\subsubsection{Inne elementy schematu – mechanizmy przetwarzania danych}
\subsubsection{Projekt mechanizmów bezpieczeństwa na poziomie bazy danych}

\subsection{Projekt aplikacji użytkownika}
\subsubsection{Architektura aplikacji i diagramy projektowe}
\subsubsection{Interfejs graficzny i struktura menu}
\subsubsection{Projekt wybranych funkcji systemu}
\subsubsection{Metoda podłączania do bazy danych – integracja z bazą danych}
\subsubsection{Projekt zabezpieczeń na poziomie aplikacji}
Na poziomie aplikacji kontrola dostępu opiera się na sesji Better Auth oraz procedurach tRPC: odczyt katalogu realizowany jest jako procedura publiczna (\texttt{publicProcedure}), natomiast wypożyczenia i rezerwacje -- jako procedury chronione (\texttt{protectedProcedure}), wymagające ważnej sesji użytkownika. Panel bibliotekarza jest dodatkowo ograniczony do roli \texttt{LIBRARIAN}.

\section{Implementacja systemu baz danych}
Implementacja i testy bazy danych w wybranym systemie zarządzania bazą danych.

\subsection{Tworzenie tabel i definiowanie ograniczeń}
\subsection{Implementacja mechanizmów przetwarzania danych}
\subsection{Implementacja uprawnień i innych zabezpieczeń}
\subsection{Testowanie bazy danych na przykładowych danych}

\section{Implementacja i testy aplikacji}
Skrócone sprawozdanie z etapu implementacji i testowania aplikacji.

\subsection{Instalacja i konfigurowanie systemu}
\subsection{Instrukcja użytkowania aplikacji}
\subsection{Testowanie opracowanych funkcji systemu}
\subsection{Omówienie wybranych rozwiązań programistycznych}
\subsubsection{Implementacja interfejsu dostępu do bazy danych}
\subsubsection{Implementacja wybranych funkcjonalności systemu}
\subsubsection{Implementacja mechanizmów bezpieczeństwa}

\section{Podsumowanie i wnioski}
% Miejsce na treść

\newpage
% --- Sekcje końcowe ---

\addcontentsline{toc}{section}{Literatura}
\begin{thebibliography}{9}
    \bibitem{ksiazka} 
    Nazwisko I., \textit{Tytuł książki}, Wydawnictwo, Rok wydania.
\end{thebibliography}

\addcontentsline{toc}{section}{Spis rysunków}
\listoffigures

\addcontentsline{toc}{section}{Spis tabel}
\listoftables

\end{document}